%% implementation.tex
%% Implementation section for the P2MAT paper.
%% Covers: package/technical details, installation, license, and limitations.

\section{Implementation}
\label{sec:implementation}

P2MAT is a cross-platform desktop application for predicting thermophysical
properties of chemical compounds from SMILES input strings.
The implementation is described below in terms of the third-party packages that
underpin the software, the technical architecture connecting them, the
installation procedures for both supported platforms, the governing license, and
the known limitations of the current release.

% ============================================================
\subsection{Packages and Technical Details}
\label{subsec:packages}

Table~\ref{tab:packages} lists the primary third-party libraries used in P2MAT.
The software requires Python 3.12 and is managed through a dedicated
\texttt{conda} environment (\texttt{qsai}) defined in \texttt{QSAI.yml}.

\begin{table}[htbp]
\centering
\caption{Key third-party packages used in P2MAT.}
\label{tab:packages}
\begin{tabular}{lll}
\hline
\textbf{Package} & \textbf{Role} & \textbf{Version channel} \\
\hline
PyQt5           & Desktop GUI framework               & pip \\
RDKit           & SMILES validation and molecular manipulation & conda-forge \\
PaDEL-Descriptor (via \texttt{padelpy}) & Molecular descriptor computation & pip \\
scikit-learn    & ML models and joblib serialization  & conda-forge \\
XGBoost         & Gradient-boosted tree regressor     & conda-forge \\
LightGBM        & Gradient-boosted tree regressor     & conda-forge \\
PyTorch         & Deep neural network (DNN) regressor & conda-forge \\
NumPy           & Numerical array operations          & conda-forge \\
Pandas          & Tabular data management             & conda-forge \\
PubChemPy       & PubChem CID $\rightarrow$ SMILES retrieval & conda-forge \\
Joblib          & Model serialization and loading     & conda-forge \\
Java (OpenJDK)  & PaDEL-Descriptor runtime            & system \\
PyInstaller     & Frozen binary packaging             & pip \\
\hline
\end{tabular}
\end{table}

\subsubsection*{Graphical User Interface}
The application window is implemented with \textbf{PyQt5}.
All blocking computation is isolated in a \texttt{QThread} subclass
(\texttt{PredictionWorker}) so that the Qt event loop remains responsive during
prediction.
The worker communicates with the main thread exclusively through four typed Qt
signals: \texttt{progress(int)}, \texttt{result(object)},
\texttt{invalid\_smiles(list)}, and \texttt{error(str)}.
The visible properties and the enabled state of each checkbox are loaded at
startup from a single \texttt{config.json} file, allowing new properties to be
toggled without recompiling or repackaging the application.

\subsubsection*{Molecular Descriptor Computation}
Molecular features are computed by \textbf{PaDEL-Descriptor}
\cite{Yap2011}, a Java-based cheminformatics tool that generates both
2D and 3D descriptors from SMILES strings.
The \texttt{padelpy} Python wrapper is used to invoke PaDEL through the
system's Java runtime, producing up to approximately 1\,240 numerical
descriptors per molecule.
To manage large query sets efficiently, SMILES are sent to PaDEL in batches of
128 molecules by default.
A divide-and-conquer fallback is applied: if a batch call fails, the batch is
recursively halved until individual molecules are attempted, ensuring that one
problematic structure does not prevent the rest from being processed.
The temporary \texttt{.smi} input and \texttt{.csv} output files created by
each PaDEL call are written to the operating-system temporary directory and
deleted immediately after reading, regardless of success or failure.

\subsubsection*{SMILES Validation}
Input SMILES strings are validated structurally using \textbf{RDKit}
\cite{Landrum2006} before any descriptor computation is attempted.
\texttt{Chem.MolFromSmiles} is called for each candidate string; a
\texttt{None} return value or an exception is treated as an invalid structure.
Invalid SMILES are returned to the caller and displayed to the user without
aborting prediction for the remaining molecules.
PubChem compound IDs are also supported: the \textbf{PubChemPy} library
retrieves the isomeric SMILES for each CID, which then passes through the same
validation path.

\subsubsection*{Property Prediction Models}
All ML models are stored as serialized \texttt{joblib} artifacts in the
\texttt{include/best\_models/} directory and are loaded on demand.
The melting-point model is a stacking ensemble comprising four base learners---
LightGBM, XGBoost, ExtraTrees, and a PyTorch deep neural network (DNN)---whose
out-of-fold predictions are combined by a Ridge meta-learner.
Boiling point, heat capacity, and heat of hydrogenation are predicted by
dedicated gradient-boosting models following the same featurisation pipeline.
H\textsubscript{2} uptake is computed rule-based by counting unsaturated
(double and triple) bonds via RDKit and converting all such bonds to single
bonds, yielding the number of moles of H\textsubscript{2} needed to
fully saturate the molecule; no ML model is involved.

\subsubsection*{Path Resolution and Packaging}
Resource paths (models, feature lists, CSS, images) are resolved through a
\texttt{helper.resource\_path} utility that detects at runtime whether the
application is running as a normal Python script or as a PyInstaller-frozen
binary (\texttt{sys.\_MEIPASS}), ensuring consistent asset access in both
deployment modes.

% ============================================================
\subsection{Installation}
\label{subsec:installation}

P2MAT is distributed as a platform-specific installer that automates all
prerequisite setup across three supported platforms.
Table~\ref{tab:install} summarises the installer, distribution format,
install location, and approximate disk footprint for each platform.
Step-by-step installation instructions for all platforms are provided in the
\texttt{README.md} file included with the software distribution.

\begin{table}[htbp]
\centering
\caption{P2MAT installation summary across supported platforms.}
\label{tab:install}
\begin{tabular}{p{3.2cm} p{3.8cm} p{4.2cm} p{1.2cm}}
\hline
\textbf{Platform} & \textbf{Installer} & \textbf{Install location} & \textbf{Disk} \\
\hline
macOS (Apple Silicon) & \texttt{install.command} (bash) & \texttt{\textasciitilde/Applications/P2MAT/} & $\sim$3.5\,GB \\
Windows (x64) & \texttt{Install-P2MAT.ps1} (PowerShell) & \texttt{\%LOCALAPPDATA\%\textbackslash P2MAT\textbackslash} & $\sim$4\,GB \\
Linux (x86\_64 / aarch64) & \texttt{install.sh} (bash) & \texttt{\textasciitilde/.local/share/P2MAT/} & $\sim$4\,GB \\
\hline
\end{tabular}
\end{table}

All three installers follow the same six-stage sequence: (1) system and
architecture check, (2) Java installation, (3) Miniconda3 installation,
(4) \texttt{qsai} conda environment creation from \texttt{QSAI.yml},
(5) application file copy, and (6) launcher or shortcut creation.
An active internet connection is required for the first install only;
subsequent launches are fully offline.
The macOS installer additionally requires \textbf{Homebrew} as the system
package manager, whereas the Linux installer auto-detects the distribution
from \texttt{/etc/os-release} and selects the appropriate package manager
(\texttt{apt}, \texttt{dnf}, \texttt{zypper}, or \texttt{pacman}) to install
OpenJDK\,17.
On Windows, Java is obtained via \texttt{winget} or a direct MSI download
if \texttt{winget} is unavailable.
An optional Inno Setup wizard installer
(\texttt{P2MAT-v1.0.0-Windows-x64-Setup.exe}) is also provided for Windows
users who prefer a guided graphical setup experience.
For advanced users who already have \texttt{conda} and Java installed, the
application can be launched directly from the repository without running the
platform installer:

\begin{verbatim}
conda env create -f QSAI.yml
conda activate qsai
cd P2MAT && python p2mat.py
\end{verbatim}

% ============================================================
\subsection{License}
\label{subsec:license}

P2MAT is released under the \textbf{BSD 3-Clause License}.

\begin{quote}
Copyright \textcopyright{} 2025, National Technology \& Engineering Solutions
of Sandia, LLC (NTESS).
Under the terms of Contract DE-NA0003525 with NTESS, the U.S.\ Government
retains certain rights in this software.
\end{quote}

Redistribution and use in source and binary forms, with or without
modification, are permitted provided that the following three conditions are
met:
(1) redistributions of source code must retain the above copyright notice, this
list of conditions, and the following disclaimer;
(2) redistributions in binary form must reproduce the above copyright notice,
this list of conditions, and the following disclaimer in the accompanying
documentation; and
(3) neither the name of NTESS nor the names of contributors may be used to
endorse or promote products derived from this software without explicit prior
written permission.

The software is provided \emph{as is}, without any express or implied warranty
of merchantability or fitness for a particular purpose.
The BSD 3-Clause License is a permissive open-source license that places
minimal restrictions on reuse while protecting the name of the originating
institution.
A copy of the full license text is included in the \texttt{LICENSE} file in
the root of the P2MAT distribution.

% ============================================================
\subsection{Limitations}
\label{subsec:limitations}

The current release of P2MAT has the following known technical and usability
limitations.

\subsubsection*{Technical Limitations}
\begin{itemize}
  \item \textbf{Platform support.}
        The macOS distribution is compiled and tested exclusively on
        Apple Silicon (M-series) hardware.
        An Intel Mac build is not provided because the melting-point DNN was
        trained using PyTorch's Metal Performance Shaders (MPS) backend; running
        the same artifact on an x86 CPU requires retraining on that hardware.
        Linux is not currently supported as a deployment target.

  \item \textbf{Java dependency.}
        PaDEL-Descriptor requires a Java runtime (OpenJDK 11 or later).
        If Java is absent or not on the system \texttt{PATH}, descriptor
        computation will fail with a timeout or subprocess error.
        The provided installers handle this automatically; manual installations
        must ensure Java is correctly configured.

  \item \textbf{PaDEL timeout and complex molecules.}
        Each molecule is subject to a per-molecule timeout of 10 seconds within
        PaDEL.
        Very large, heavily branched, or structurally unusual molecules may
        exhaust this budget, causing them to be silently skipped after the
        divide-and-conquer fallback is exhausted.
        Processing such molecules individually, or increasing the timeout
        constant in \texttt{utils/descriptor.py}, may mitigate this issue.

  \item \textbf{Applicability domain.}
        All ML models (melting point, boiling point, heat capacity, heat of
        hydrogenation) were trained on finite curated datasets of predominantly
        organic molecules.
        Predictions for structures that lie outside the chemical space covered
        by the training data---such as inorganic salts, organometallics, or
        highly unusual functional groups---may be unreliable.
        No applicability-domain filter is currently applied automatically,
        so it is the user's responsibility to interpret results for atypical
        inputs with caution.

  \item \textbf{Internet requirement.}
        An active internet connection is required during the initial
        installation to download Miniconda, package dependencies, and OpenJDK.
        Offline installation is not supported.
        PubChem CID lookup (\texttt{prediction\_from\_pid}) also requires
        network access at prediction time.
\end{itemize}

\subsubsection*{Usability Limitations}
\begin{itemize}
  \item \textbf{Input format.}
        P2MAT accepts only SMILES strings or PubChem compound IDs as input.
        Other common molecular representations (InChI, MOL/SDF files, CAS
        registry numbers) are not natively supported and must be converted to
        SMILES externally before use.

  \item \textbf{Sequential per-molecule processing.}
        Molecules are processed one at a time in the background worker.
        Although the UI remains responsive, the total wall-clock time scales
        linearly with the number of input molecules because descriptor
        computation involves launching the Java-based PaDEL subprocess for each
        SMILES independently (after internal batching).
        Large batches of hundreds of molecules may require several minutes.

  \item \textbf{Single-session results.}
        Prediction results are held in memory for the duration of the session.
        There is no built-in history or session persistence; results that are
        not exported to CSV before the window is closed are lost.

  \item \textbf{H\textsubscript{2} uptake model scope.}
        The H\textsubscript{2} uptake feature is a rule-based estimate based
        solely on unsaturation count (double and triple bonds) and does not
        account for steric accessibility, partial hydrogenation, or
        thermodynamic feasibility.
        It should be treated as a stoichiometric upper bound rather than a
        kinetic or thermodynamic prediction.

  \item \textbf{Minimum one property constraint.}
        The GUI enforces that at least one property checkbox remains checked at
        all times.
        If all checkboxes are unchecked, melting point is automatically
        re-enabled as a fallback.
        This prevents zero-property prediction runs but may be unexpected for
        users who intend only to validate SMILES input without running a model.
\end{itemize}
